\subsection*{Moment transport from the Appendix-A reference radius}
The source's reference region (A.7) has
$U=4\eta$, $E=P_*f(\eta)(X/X_R)^{1/10}$,
$f=(1+\eta^2)^{-1}$, and $H=\sqrt{2X}E$.
The cumulative moments are $M=\int_0^XU\,dx$,
$J=\int_0^XUH\,dx$, and
$S=\int_0^X(U^2-E^2/2)\,dx$, as in (4.15).

\begin{proposition}[Exact source moment anchor]
At the first matching radius,
\begin{equation}\label{eq:source-ref-MJS}
\begin{aligned}
M(X_R)&=4\eta X_R,\\
J(X_R)&=\frac52\eta X_R\sqrt{2X_R}\,P_*f,\\
S(X_R)&=X_R\left(16\eta^2-\frac5{12}P_*^2f^2\right).
\end{aligned}
\end{equation}
\end{proposition}
\begin{proof}
Integrate the source's reference power laws. The three needed
integrals over $0\le x=X/X_R\le1$ are
$\int_0^1dx=1$,
$\int_0^1x^{3/5}dx=5/8$ and
$\int_0^1x^{1/5}dx=5/6$.
\end{proof}

Put $y=\log(X/X_R)$,
$a=U/E$, $\ell=\partial_y\log(\sqrt XE)$,
and define the scale-invariant moments
\[
m=\frac{M}{XE},\qquad j=\frac{J}{XHE},\qquad
s=\frac{S}{XE^2},\qquad
d=\frac1{2XE^2}\int_X^\infty E(x,\eta)^2\,dx.
\]
\begin{theorem}[Exact four-moment transport]\label{thm:four-moment-transport}
On every source stage,
\begin{equation}\label{eq:four-moment-transport}
\begin{aligned}
m'&=a-(\ell+\tfrac12)m,&
j'&=a-(2\ell+\tfrac12)j,\\
s'&=a^2-\tfrac12-2\ell s,&
d'&=-\tfrac12-2\ell d.
\end{aligned}
\end{equation}
The forward initial values at $X_R$ are
$m=4\eta/(P_*f)$,
$j=5\eta/(2P_*f)$ and
$s=16\eta^2/(P_*^2f^2)-5/12$.
At the start $X_T$ of the source's pure power-law tail,
$d(X_T)=1/(4h)$.
If $X_p$ is the pulse start and
$X_{\mathrm{end}}=X_pe^{13/\lambda}$ its end, the
amplitude-independent remainder from Proposition~\ref{prop:source-exact-A19}
obeys the exact identity
\begin{equation}\label{eq:source-bias-transport}
P_{\mathrm{pre}}-D_{\mathrm{post}}
=\lambda\bigl(s(X_p)-e^{-26}d(X_{\mathrm{end}})\bigr).
\end{equation}
\end{theorem}
\begin{proof}
The logarithmic derivatives of $XE$, $XHE$, and $XE^2$
are $\ell+1/2$, $2\ell+1/2$, and $2\ell$,
respectively. Use $dM/dy=XU$,
$dJ/dy=XUH$,
$dS/dy=X(U^2-E^2/2)$, and
$d(\frac12\int_X^\infty E^2dx)/dy=-XE^2/2$
to prove the four transport equations.
For $X\ge X_T$, $E=E_T(X/X_T)^{-1/2-h}$;
hence $\frac12\int_{X_T}^\infty E^2dx=X_TE_T^2/(4h)$.
On the pulse,
$XE^2/(X_pe_b^2f^2)=e^{-2\lambda\log(X/X_p)}$;
at $X_{\mathrm{end}}$ this is $e^{-26}$. The bias
identity follows directly from its definitions.
\end{proof}
The first three equations admit forward integrating-factor
formulas; the last is integrated backward from $d(X_T)=1/(4h)$.
An interval calculation need only enclose $a(y,\eta)$ and
$\ell(y,\eta)$ on the \emph{actual} finitely many Appendix-A
stages. It is not necessary to introduce an artificial
semi-infinite pre-pulse profile or an approximate terminal tail.

\subsection*{Correction to the prescribed moment-bump estimate}
Appendix A fixes the positive correction bumps as rescaled copies
of $\sigma'$ from (A.5), on intervals of width $3/10$.
The former use of an alternative normalized plateau bump with
$\int\beta^2\le5$ was not a certificate for this actual
source shape. For the prescribed bump,
\[
\beta(y)=\frac{10}{3}\sigma'\left(
 \frac{10}{3}(y-y_c)+\frac12\right),\qquad
\sigma'(t)=\sigma(t)(1-\sigma(t))
\left(\frac2{t^3}+\frac2{(1-t)^3}\right).
\]
\begin{proposition}[Prescribed source-bump norm and corrected error]
\label{prop:true-source-bump-L2}
The actual normalized smooth-step-derivative bump satisfies
\begin{equation}\label{eq:source-bump-correct}
\frac{27}{5}<B_2:=\int\beta(y)^2\,dy\le\frac{640}{3}.
\end{equation}
If the actual transported source moments satisfy
$\max(|m(X_p)|,|j(X_p)|)\le1$, then for
$0<\lambda\le1/20$, $|A|\le1.011$ the
complete amplitude remainder satisfies
\[
0\le\mathcal E-(P_{\mathrm{pre}}-D_{\mathrm{post}})
<5\cdot10^{-17},\qquad
|\partial_A\mathcal E|<10^{-16}.
\]
\end{proposition}
\begin{proof}
The symmetry $\sigma(1-t)=1-\sigma(t)$ gives
$\sigma(1/4)=1/(1+e^{128/9})<1/20$,
and $\sigma(3/4)>19/20$.
Cauchy--Schwarz on $[1/4,3/4]$ yields
$\int_0^1\sigma'^2>2(9/10)^2$.
On $[1/4,3/4]$ one has $\sigma'\le64$.
For $0<t\le1/4$,
$\sigma'\le e^{2-16}(128+128/27)<1$ using
monotonicity of $t^{-3}e^{-1/t^2}$.
Symmetry handles the remaining quarter. Hence
$\int_0^1\sigma'^2\le64\int_0^1\sigma'=64$,
and scaling by $10/3$ proves the norm bounds.

Insert $B_2\le640/3$ into the already proved
positive remainder inequalities. With $t=1/\lambda\ge20$,
$D_0+1.011D_1\le(57605/2000)e^{6/5-4t/5}$.
The rigorously verified exponential bounds
$e^{-74/5}<4\cdot10^{-7}$,
$e^{-5137/200}<10^{-11}$ and monotonicity of
$t e^{-8t/5+63/(10t)}$ give exact rational upper
bounds $4.531\cdot10^{-17}$ and
$8.652\cdot10^{-17}$ for the two errors.
\end{proof}
A fully source-admissible amplitude interval therefore requires
\[
\max(|m(X_p)|,|j(X_p)|)\le1,\qquad
\lambda\sup_{|\eta|\le1}
|s(X_p,\eta)-e^{-26}d(X_{\mathrm{end}},\eta)|
\le\frac1{2500}-5\cdot10^{-17}.
\]
Neither condition has yet been numerically certified for one
frozen, moment-matched exterior schedule. The corrected
remainder and transport formulas establish no unforced
Navier--Stokes solution or physical residual theorem.
